% ============================================================================
% SEC06.TEX - Section 3.6: Coroutine
% ============================================================================

\section{Coroutine}

\begin{learningoutcomes}
\item Memahami konsep coroutine
\item Menggunakan IEnumerator dan yield
\item Menerapkan coroutine untuk delay dan timing
\item Menggunakan coroutine untuk async operations
\end{learningoutcomes}

\subsection{Pengertian Coroutine}

\begin{definisi}[Coroutine]
Coroutine adalah fungsi yang dapat menjeda eksekusinya sendiri dan melanjutkan di frame berikutnya, memungkinkan operasi yang berjalan selama beberapa frame \cite{UnityScriptingAPI2024}.
\end{definisi}

\subsection{Contoh Implementasi}

\begin{lstlisting}[language={[Sharp]C}]
using System.Collections;
using UnityEngine;

public class CoroutineExample : MonoBehaviour
{
    void Start()
    {
        StartCoroutine(Countdown());
    }
    
    IEnumerator Countdown()
    {
        for (int i = 5; i > 0; i--)
        {
            Debug.Log(i);
            yield return new WaitForSeconds(1f);
        }
        Debug.Log("Go!");
    }
    
    IEnumerator FadeOut()
    {
        float alpha = 1f;
        while (alpha > 0)
        {
            alpha -= Time.deltaTime;
            // Update alpha
            yield return null; // Wait one frame
        }
    }
}
\end{lstlisting}

\begin{catatan}
Coroutine sangat berguna untuk operasi yang perlu berjalan selama beberapa frame seperti animasi, delay, atau loading.
\end{catatan}
